% =====================================================================
%  Seance 4 -- D. TP : prevoir et decider (frames 22--27)
% =====================================================================
\section{TP « Prévoir et décider »}

\begin{frame}[fragile]{TP B1 --- la série horaire du parc (30 min)}
\begin{codebox}[B1. Agréger, décomposer, tester]
\begin{lstlisting}
import pandas as pd
from statsmodels.tsa.seasonal import seasonal_decompose
from statsmodels.tsa.stattools import adfuller

m = pd.read_csv("data/nosys_metriques.csv",
                parse_dates=["horodatage"])
s = m.set_index("horodatage")["cpu_pct"].resample("1h").mean()
print(len(s))                     # 720

dec = seasonal_decompose(s, model="additive", period=24)
dec.plot()                        # tendance + saison + residus

stat, p, *_ = adfuller(s)         # 0.075 : on differencie
stat, p, *_ = adfuller((s - s.shift(24)).dropna())  # ~0
\end{lstlisting}
\end{codebox}
\begin{abox}[Livrable B1]
Un graphique à 4 panneaux : série, décomposition, ACF, PACF
(\py{plot_acf}, \py{plot_pacf} de \py{statsmodels}) + la conclusion
ADF en une phrase. Vos chiffres doivent retrouver ceux du cours :
$p_{\text{brut}} = 0{,}075$, saison à 24 h.
\end{abox}
\end{frame}

\begin{frame}[fragile]{TP B2 --- SARIMA : ajuster, prévoir, mesurer (40 min)}
\begin{codebox}[B2. Le protocole du cours, exécuté par vous]
\begin{lstlisting}
from statsmodels.tsa.statespace.sarimax import SARIMAX
import numpy as np

train, test = s.iloc[:576], s.iloc[576:]
t = np.arange(len(s))             # harmoniques hebdo (exogenes)
X = np.column_stack(
    [f(t*2*np.pi*j/168) for j in (1,2,3)
     for f in (np.sin, np.cos)])

mod = SARIMAX(train, exog=X[:576],
              order=(2,0,1), seasonal_order=(1,0,1,24))
res = mod.fit(disp=False)
fc = res.get_forecast(144, exog=X[576:])
pred, ci = fc.predicted_mean, fc.conf_int()

mae = np.mean(np.abs(pred - test))     # ~3.0
\end{lstlisting}
\end{codebox}
\begin{wbox}[Deux contrôles obligatoires]
\begin{enumerate}
\item MAE \textbf{mesurée} sur les 144 heures de test, jamais sur le
train --- et comparée à deux bases naïves (persistance 24 h : 5,06 ;
semaine précédente : 5,23).
\item Toujours tracer \textbf{l'intervalle} : c'est lui qui porte la
marge de capacité du cours 3.
\end{enumerate}
\end{wbox}
\end{frame}

\begin{frame}{B2 --- ce que vous devez observer}
\begin{center}
\includegraphics[width=0.8\textwidth]{f26_sarima_prevision.png}
\end{center}
\vspace{1mm}
\begin{columns}[t]
\begin{column}{0.5\textwidth}
\begin{solbox}[Résultat de référence]
\small
MAE test \textbf{3,01 pts CPU} : \textbf{$-40{,}5$ \%} d'erreur par
rapport à la persistance. L'intervalle 95 \% couvre l'essentiel des
pics --- sauf les incidents, invisibles pour un modèle de niveau.
\end{solbox}
\end{column}
\begin{column}{0.46\textwidth}
\begin{qbox}[Question d'analyse]
\small
Pourquoi l'intervalle s'élargit-il avec l'horizon ? (réponse :
l'incertitude de prévision s'accumule --- au-delà de quelques jours,
prévoir et moyenne hebdomadaire se valent.)
\end{qbox}
\end{column}
\end{columns}
\end{frame}

\begin{frame}[fragile]{TP B3 --- modèle de panne, split temporel strict (30 min)}
\begin{codebox}[B3. Le split qui respecte le temps]
\begin{lstlisting}
from sklearn.ensemble import GradientBoostingClassifier
from sklearn.metrics import confusion_matrix

f = pd.read_csv("data/nosys_features.csv",
                parse_dates=["fenetre"])
feats = ["cpu_moy","cpu_max","ram_moy","lat_moy","lat_p99",
         "perte_moy","temp_moy","conso_moy","cpu_delta",
         "cpu_amplitude","h_sin","h_cos"]
f = f.dropna(subset=feats+["incident_30min"])

cut = pd.Timestamp("2026-09-22")        # JAMAIS de shuffle
tr, te = f[f.fenetre < cut], f[f.fenetre >= cut]
gbc = GradientBoostingClassifier(
    n_estimators=300, max_depth=4,
    learning_rate=0.06, random_state=5534)
gbc.fit(tr[feats], tr.incident_30min)
print(confusion_matrix(te.incident_30min,
      gbc.predict(te[feats])))
# -> [[10290, 14], [24, 40]] : rappel 62,5 %
\end{lstlisting}
\end{codebox}
\begin{abox}[Vérification anti-fuite]
Refaire le modèle avec \py{alarme} en entrée : 40 fenêtres d'alarme
donnent 40 incidents « prédits » à 100 \%. C'est le signe de la fuite
--- la variable répète l'étiquette au lieu de la précéder.
\end{abox}
\end{frame}

\begin{frame}{TP B4 --- le budget d'alertes (20 min)}
\begin{columns}[t]
\begin{column}{0.52\textwidth}
\begin{block}{Principe}
Le rappel seul ne décide rien : chaque fausse alerte coûte du temps
d'astreinte. On trace le \textbf{rappel en fonction du nombre
d'alertes par jour} (classer les fenêtres de test par score
décroissant, regarder les $k$ premières chaque jour) et on choisit le
\textbf{point d'opération}.
\end{block}
\begin{solbox}[Référence du dossier]
\textbf{5 alertes/jour} $\Rightarrow$ rappel \textbf{62,5 \%}
(40 pannes sur 64 rattrapées). Au-delà, la courbe plafonne : les
24 incidents restants n'ont aucun signal métrique --- il faudra les
journaux (séance 5).
\end{solbox}
\end{column}
\begin{column}{0.44\textwidth}
\begin{qbox}[À débattre en binôme]
Votre équipe accepte 2 alertes/jour maximum. Quel rappel obtenez-vous ?
Quel compromis proposeriez-vous au responsable d'exploitation, chiffré
en incidents manqués et en heures d'astreinte ?
\end{qbox}
\end{column}
\end{columns}
\end{frame}

\begin{frame}{TP B5 --- livrable et barème}
\begin{columns}[t]
\begin{column}{0.52\textwidth}
\begin{abox}[À rendre]
\begin{enumerate}
\item \textbf{Graphique prévision/réalisé} avec intervalle (B2).
\item \textbf{Tableau rappel / alertes par jour} (B4) avec le point
d'opération choisi et justifié.
\item \textbf{Date de saturation} du stockage, avec hypothèses
explicites (cours 3).
\item \textbf{Un paragraphe} : « que décide l'exploitant lundi matin ? »
--- scale, marge, rétention, seuils.
\end{enumerate}
\end{abox}
\end{column}
\begin{column}{0.44\textwidth}
\begin{kbox}[Barème (20 pts)]
\begin{itemize}
\item B1 série, décomposition, ADF --- 4 pts
\item B2 SARIMA ajusté + MAE honnête --- \textbf{5 pts}
\item B3 modèle sans fuite + split strict --- \textbf{5 pts}
\item B4 courbe et choix du point --- 3 pts
\item B5 livrable et décision argumentée --- 3 pts
\end{itemize}
\end{kbox}
\end{column}
\end{columns}
\end{frame}
